\documentclass[11pt]{article}
\usepackage[margin=2.54cm]{geometry}
\usepackage[dvipsnames,usenames]{color}
\usepackage{amsfonts,amsmath,amssymb,amsthm,mathtools}
\usepackage{paralist}
\usepackage{algorithmic,algorithm}
\usepackage{bm}
\usepackage{xspace}
\usepackage[pagebackref,letterpaper=true,colorlinks=true,pdfpagemode=none,urlcolor=blue,linkcolor=blue,citecolor=BrickRed,pdfstartview=FitH]{hyperref}
\usepackage{xspace,prettyref}
\usepackage{color}
%\usepackage{MnSymbol}
%% For hyperlinks. Should always be the last package.
%\usepackage[colorlinks,urlcolor=blue,citecolor=blue,linkcolor=blue]{hyperref}
\let\pref=\prettyref
\newcommand{\savehyperref}[2]{\texorpdfstring{\hyperref[#1]{#2}}{#2}}
\newcommand{\comment}[1]{ {\color{BrickRed} \footnotesize[#1]}\marginpar{\footnotesize\textbf{\color{red} To Do!}}}
\newcommand{\ignore}[1]{}

\newtheorem{theorem}{Theorem}
\newtheorem{lemmata}[theorem]{Lemmata}
\newtheorem{lemma}[theorem]{Lemma}
\newtheorem{claim}[theorem]{Claim}
\newtheorem{subclaim}{Subclaim}
\newtheorem{proposition}[theorem]{Proposition}
\newtheorem{corollary}[theorem]{Corollary}
\newtheorem{fact}[theorem]{Fact}
\newtheorem{conjecture}[theorem]{Conjecture}
\newtheorem{question}[theorem]{Question}
\newtheorem{example}[theorem]{Example}
\newtheorem{definition}{Definition}
\theoremstyle{definition}
\newtheorem{remark}{Remark}
\newtheorem{observation}{Observation}

\def\R{\RR}
\def\RR{\mathbb R}
\def\ZZ{\mathbb Z}
\def\eqdef{~\triangleq~}
\def\eps{\varepsilon}
\def\bar{\overline}
\def\floor#1{\lfloor {#1} \rfloor}
\def\ceil#1{\lceil {#1} \rceil}
\def\script#1{\mathcal{#1}}

\newenvironment{proofof}[1]{\smallskip\noindent{\bf Proof of #1:}}%
        {\hspace*{\fill}$\Box$\par}

%% Hypercube related macros

\def\B{{\mathsf D}}
\def\Hyp{{\sf Hyp}}
\def\D{{\sf K}}
\def\G{{\mathsf G}}
\def\d{{\sf d}}
\def\a{{\mathfrak a}}
\def\b{{\mathfrak b}}
\def\c{{\mathfrak c}}
\def\db{{\mathfrak d}}
\def\e{{\mathfrak e}}
\def\f{{\mathfrak f}}
\def\h{{\mathfrak h}}
\def\g{{\mathfrak g}}

\newcommand{\PTIME}{\mathbb{P}}
\newcommand{\NP}{\mathbb{NP}}
\newcommand{\coNP}{\textrm{co-}\mathbb{NP}}


\newcommand{\hu}{\widehat{u}}

\newcommand{\hcd}{\mathsf{hcd}}
\newcommand{\VG}{VG}
\newcommand{\md}[1]{\ (\operatorname{mod} #1)}

\newcommand{\cupdot}{\sqcup}

%% Calligraphic letters

\newcommand{\cA}{{\cal A}}
\newcommand{\cB}{{\cal B}}
\newcommand{\cC}{{\cal C}}
\newcommand{\cD}{{\cal D}}
\newcommand{\cE}{{\cal E}}
\newcommand{\cF}{{\cal F}}
\newcommand{\cG}{{\cal G}}
\newcommand{\cH}{{\cal H}}
\newcommand{\cI}{{\cal I}}
\newcommand{\cJ}{{\cal J}}
\newcommand{\cL}{{\cal L}}
\newcommand{\cM}{{\cal M}}
\newcommand{\cP}{{\cal P}}
\newcommand{\cR}{{\cal R}}
\newcommand{\cS}{{\cal S}}
\newcommand{\cT}{{\cal T}}
\newcommand{\cU}{{\cal U}}
\newcommand{\cV}{{\cal V}}

%% Bold letters

\newcommand{\bS}{{\bf S}}

\newcommand{\NN}{\mathbb{N}}

%% Hyper-linked References
\newcommand{\Sec}[1]{\hyperref[sec:#1]{\S\ref*{sec:#1}}} %section
\newcommand{\Eqn}[1]{\hyperref[eq:#1]{(\ref*{eq:#1})}} %equation
\newcommand{\Fig}[1]{\hyperref[fig:#1]{Fig.\,\ref*{fig:#1}}} %figure
\newcommand{\Tab}[1]{\hyperref[tab:#1]{Tab.\,\ref*{tab:#1}}} %table
\newcommand{\Thm}[1]{\hyperref[thm:#1]{Theorem\,\ref*{thm:#1}}} %theorem
\newcommand{\Lem}[1]{\hyperref[lem:#1]{Lemma\,\ref*{lem:#1}}} %lemma
\newcommand{\Prop}[1]{\hyperref[prop:#1]{Prop.~\ref*{prop:#1}}} %property
\newcommand{\Cor}[1]{\hyperref[cor:#1]{Corollary~\ref*{cor:#1}}} %corollary
\newcommand{\Def}[1]{\hyperref[def:#1]{Definition~\ref*{def:#1}}} %definition
\newcommand{\Alg}[1]{\hyperref[alg:#1]{Alg.~\ref*{alg:#1}}} %algorithm
\newcommand{\Ex}[1]{\hyperref[ex:#1]{Ex.~\ref*{ex:#1}}} %example
\newcommand{\Clm}[1]{\hyperref[clm:#1]{Claim~\ref*{clm:#1}}} %example


\def\yes{{\sf YES }}
\def\no{{\sf NO }}
\def\Lip{{\sf Lip}}
\def\poly{{\sf poly}}

\newcommand{\ab}{\hbox{ab}}
\newcommand{\abs}[1]{{|#1|}}
\newcommand{\hd}[2]{{||#1 - #2||_1}}
\newcommand{\proj}{{\mathsf{proj}}}
\newcommand{\drop}{\hbox{drop}}
\newcommand{\dist}[2]{\eps_{#1,#2}}
\newcommand{\hM}{\widehat{M}}
\newcommand{\hX}{\widehat{X}}

\newcommand{\flo}[1]{\lfloor #1 \rfloor}

\newcommand{\Sesh}[1]{{\color{red} Sesh says: #1}}

\title{CSE104: Exercises\\Diagonalization\\{\small \copyright C. Seshadhri, 2025}}

\date{}

\begin{document}

\maketitle

Please cite any collaborations
or online sources. All assignments must be done in Latex and submitted as pdfs.

\section{Easier problems}

Try to make some attempts by yourself, before collaborating with others. Ideally, do not use GenAI for these
problems. The answers are not that long. We can use any previous statement to prove subsequent ones.

\begin{enumerate}
    \item Prove that the smallest infinite cardinal number is $\aleph_0$. (In other words: consider a set $S$ that is infinite, and has an injection (a 1-1 mapping) into $\NN$. Prove that there exists a bijection between $S$ and $\NN$.) You can use this statement in your later proofs; things become much easier.
    \item For two infinite sets $A, B$, we say that $|A| \leq |B|$ if there exists a 1-1 map (injection) between $A$ and $B$. 
Prove: if $|A| \leq |B$ and $|B| \leq |C|$, then $|A| \leq |C|$.
    \item Consider any finite alphabet $\Sigma$, and consider $\Sigma^*$, the set of finite strings over $\Sigma$. Prove that
        there is a bijection from $\Sigma^* \to \NN$. (Hint: think encodings.
        And use this fact in the next problems.)
    \item Prove that $|\NN^2| = |\NN|$. So show a bijection between the set of pairs $(x,y)$ of natural numbers and the set of natural numbers.
    \item Prove that $|\mathbb{Q}| = \aleph_0$. This means that the cardinality of rational numbers is the same as the cardinality of natural numbers.
        (Hint: feel free to use the statements above, and you will see a short proof.)
\end{enumerate}


\section{ChatGPT questions} 

Put this question ChatGPT and get the answer. You are allowed to discuss further. Please put a link to the response and your conversation. 
Given the response, you must first say whether: (i) the response is correct, (ii) the response is false, or (iii) you are really not sure.
In the first case, give a succinct explanation of the proof. In the second case, explain why it is false.
In the third case, explain your confusion. You do not need to write too much.

\begin{enumerate}
    \item Give a bijection from $\{x | x > 1, x \in \RR\}$ and the interval $(0,1)$.
\end{enumerate}


\section{More challenging exercises}

\begin{enumerate}
    \item Prove that $\bigcup_{d \in \NN} \RR^d$ and $\RR$ have the same cardinality. 
    \item The Cantor-Schr\"{o}der-Bernstein theorem states that if $|A| \leq |B|$ and $|B| \leq |A|$, then $|A| = |B|$. This is a fundamental theorem
        that shows that the infinities form an ordering.

        We will prove this theorem in this exercise. By assumption, there is an injection $f: A \to B$ and an injection $g: B \to A$. Construct a colored bipartite graph between
        $A$ and $B$ as follows. For each pair $(a, f(a))$, insert the red edge $(a, f(a))$. For each pair $(b, g(b))$, insert
        the blue edge $(g(b), b)$. (Note that $g(b) \in A$, so we flipped the order to describe the edge.) An \emph{alternating path} is 
        a finite path of edges that keeps switching colors: so either red, blue, red, blue, $\ldots$, or blue, red, blue, red, $\ldots$.
        \begin{enumerate}
            \item Definite relation $a \to b$ if there is an alternating path from $a$ to $b$. Prove that this relation is an equivalence relation.
            \item Prove that $A \cup B$ can be partitioned into set of: (i) alternating cycles, (ii) finite alternating paths, (iii)
                semi-infinite alternating paths, or (iv) doubly infinite alternating paths. (A semi-infinite alternating path has one endpoint;
                a doubly infininte alternating path has no endpoint.)
            \item Construct a bijection between $A$ and $B$, going over the four cases above.
        \end{enumerate}
\end{enumerate}
\end{document}
